\section{Cyclic trisections and the ternary trace code}

\subsection{Trisection shells in $E_8$}

Let $S$ be the split $E_8$ rational elliptic surface, with zero section $O$
and fibre class $F$.  Its N\'eron--Severi lattice decomposes as
\[
\NS(S)=\langle O,F\rangle\oplus E_8(-1),
\qquad
O^2=-1,\quad O\cdot F=1,\quad F^2=0.
\]
Let $D$ be a genus-one trisection and put
\[
m=D\cdot O\ge0.
\]
Writing
\[
D=3O+(m+3)F+v,
\qquad v\in E_8(-1),
\]
and using $K_S=-F$, adjunction gives $D^2=3$.  Therefore
\[
\boxed{\|v\|_{E_8}^2=6(m+1).}
\]
The theta series
\[
\Theta_{E_8}(q)=1+240\sum_{n\ge1}\sigma_3(n)q^n
\]
shows that the number of such lattice classes is
\[
\boxed{240\sigma_3(3(m+1)).}
\]
For $m=0$ this gives exactly $6720$ minimal trisection classes.

\subsection{Universal projected orbit geometry}

Assume the normalization of $D$ is a cyclic cubic genus-one cover with simple
total ramification over three smooth fibres.  On the pullback, let
\[
P_0=P,\qquad P_1=\sigma P,\qquad P_2=\sigma^2P.
\]
Distinct conjugates meet transversely at the three ramification points and
nowhere else, so
\[
P_i\cdot P_j=3\qquad(i\ne j).
\]
The pulled-back surface has $\chi=3$ and no reducible fibres.  Shioda's
formula yields
\[
\langle P_i,P_i\rangle=6+2m,
\qquad
\langle P_i,P_j\rangle=2m\quad(i\ne j).
\]
Let
\[
T=P_0+P_1+P_2,
\qquad
q_i=P_i-\frac{T}{3}.
\]
Then
\[
\langle T,T\rangle=18(m+1)
\]
and, independently of $m$,
\[
\boxed{
\langle q_i,q_i\rangle=4,
\qquad
\langle q_i,q_j\rangle=-2\quad(i\ne j).
}
\]
Thus every cyclic trisection orbit contributes a projected plane
\[
A_2(2),\qquad
\operatorname{Gram}=
\begin{pmatrix}4&-2\\-2&4\end{pmatrix}.
\]
The integral difference lattice generated by $P_0-P_1$ and $P_1-P_2$ is
\[
A_2(6),\qquad
\operatorname{Gram}=
\begin{pmatrix}12&-6\\-6&12\end{pmatrix}.
\]

This corrects an attractive but impossible earlier fingerprint.  Because the
surface has $\chi=3$ and no reducible fibres, every nonzero section has height
at least six.  The full Mordell--Weil lattice cannot contain roots of norm two
or four; norm four appears only after orthogonal projection away from the
invariant $E_8(3)$ lattice.

\subsection{The mod-$3$ trace code}

The trace $T$ is the pullback of the $E_8$ vector $v$.  For two orbit sections
$P,Q$ with traces $v,w$, Galois invariance gives
\[
\boxed{
\langle q_P,q_Q\rangle
=
\langle P,Q\rangle-\frac{\langle v,w\rangle_{E_8}}{3}.
}
\]
Consequently, an integral projected packet requires the trace residues to be
mutually orthogonal in $E_8/3E_8$.  Every trisection trace is isotropic modulo
three because its norm is divisible by six.

The finite quadratic space $E_8/3E_8$ is split of dimension eight and Witt
index four.  Exact enumeration gives:
\begin{itemize}[leftmargin=*]
\item $2240$ nonzero isotropic vectors;
\item exactly three norm-six $E_8$ lifts of each nonzero isotropic residue;
\item $2240$ maximal totally isotropic four-spaces;
\item $80$ nonzero residues and hence $240$ minimal trisection classes in each maximal four-space;
\item $80$ maximal four-spaces through each nonzero isotropic residue.
\end{itemize}
The projective-quadric counts are
\[
(3^3+1)(3^4-1)=2240
\]
for nonzero isotropic vectors and
\[
2\prod_{i=1}^{3}(3^i+1)=2240
\]
for maximal isotropic four-spaces.  The three-to-one lifting statement is
verified by the dependency-free exact enumeration shipped with the paper.

\subsection{Determinant identity}

Let $M$ be the full Mordell--Weil lattice of the cubic pullback and let
\[
L=M^{\langle\sigma\rangle}\cong E_8(3).
\]
Let $K$ be the image of $M$ under orthogonal projection to the trace-zero
rational subspace.  If $K$ is integral, the projection exact sequence and an
orthogonal covolume calculation give
\[
\boxed{\det M=\det L\,\det K=3^8\det K.}
\]
This determinant is a powerful adversarial check: eleven apparent orbit
planes must be tested not only for rational independence but also for ternary
gluing and saturation.

\subsection{Finite search reduction}

The minimal integral track is no longer an undifferentiated search through all
$6720$ norm-six classes.  It is the finite incidence problem
\[
\boxed{
2240\text{ trace codes}\times240\text{ minimal classes per code}.
}
\]
For each code, one computes the four-dimensional systems $|D_v|$, retains
members whose cubic discriminant is a square and whose branch divisor consists
of three smooth fibres, canonicalizes the resulting covers, and groups all
trisections defining the same cover.  Eleven Hermitian-independent orbit
planes over one positive-rank cover would contribute twenty-two new
directions; together with $E_8(3)$ they would force rank thirty.
